\section{Pinning the gradings}\label{sec:G}

We now combine the character supply of Section~\ref{sec:D} with the
optimization and transfer results of Sections~\ref{sec:E} and~\ref{sec:F}.
The discrepancy gap in this section is always a gap for primitive divisors
on the actual quotient germs.  No such gap is imposed on the auxiliary
central cones.

The local phase assertion of \cite[Theorem~3.1]{OpenAILocalPhase2026}
constructs centred divisorial valuations with bounded homogeneous discrepancy
and cyclic-character congruences, after passage to a subsequence, under its
hypotheses on Gorenstein klt germs of fixed dimension $n\ge2$.
It does not guarantee that these valuations lie on the prescribed $h$-lc
slice $A(T)=T(h)=1$ below, so it does not replace the pinning obstruction.

\begin{theorem}[Pinning obstruction]\label{thm:G-pinning}
Fix an integer $n>0$ and $\epsilon>0$.  There is no sequence of the following
data satisfying all the conditions below.  For each index $i$, let
$x_i\in S_i$ be an $n$-dimensional complex klt germ with an action of a
finite cyclic group $\Gamma_i$ fixing $x_i$, such that the quotient map is
quasi-\'etale and $S_i/\Gamma_i$ is $\epsilon$-lc.  Suppose $K_{S_i}$ has a
nowhere-vanishing local eigen-generator $\omega_i$, and let $h_i$ be a
nonzero regular eigenfunction with the same character $\chi_i$ as
$\omega_i$.  The character need not be faithful.  Assume that
$(S_i,\divisor(h_i))$ is lc and admits a centred quasi-monomial valuation
$T_i$ with
\[
 A_{S_i}(T_i)=T_i(h_i)=1.
\]
The excluded additional condition is: for every fixed integer $m>0$,
there is an index $i(m)$ such that, for all $i\ge i(m)$ and every regular
semi-invariant $s$ of character $\chi_i^m$,
\begin{equation}\label{eq:G-input}
                         T_i(s)\ge m.
\end{equation}
\end{theorem}

All constants in the proof are uniform along the sequence under
consideration; they may depend on the dimension and on the finitely many
previous stages of the construction.  An assertion made eventually also
holds after any further subsequence.  We suppress $i$ whenever this does
not obscure a limit.  We may successively pass to subsequences, since the
number of stages and of subsidiary row additions will be bounded.

\subsection{Costs and the induction data}

An invariant rational function $Y$ has \emph{cost at most $m$} if
$h^mY$ is regular, where $m$ is a nonnegative integer.  We increase costs
to positive integers and to common denominators when convenient.  Costs
are subadditive under multiplication.  The same terminology applies on
a graded cone, with $h$ replaced by its initial symbol.  Every homogeneous
ratio on such a cone of cost $m$ lifts to an actual invariant ratio of the
same cost: lift its regular numerator, and project to the character
$\chi^m$ by linear reductivity of the finite group.

We keep the test functions and accumulated variables on the original germ,
while replacing the working valuation and its cone.  Algebraic independence
of the accumulated symbols is the quantity that will force termination.
Their weights may nevertheless become rationally proportional: suitable
ratios of monomials then have weight zero and retain information in the
angular field.  Pinning the weights to a line therefore need not reduce
the transcendence degree of the accumulated symbols.

The following data are retained at the beginning of each stage.
\begin{enumerate}[label=\textup{(I\arabic*)}]
\item\label{it:G-tests}
There are finitely many actual invariant test ratios $f_j$, each of cost
$O(L_-)$, where initially $L_-=1$.  The common actual slice
\begin{equation}\label{eq:G-slice}
 \mathcal Q=\{q:A_S(q)=q(h)=1,\ q(f_j)\ge0\text{ for every }j\}
\end{equation}
contains a centred valuation, and every test has value exactly zero
throughout this slice.

\item\label{it:G-pairs}
There are accumulated pairs $Y_j^+,Y_j^-$ of invariant ratios of cost
$O(L_-)$, with opposite values on $\mathcal Q$.  We write
$y_j(q)=q(Y_j^+)=-q(Y_j^-)$.  The products $Y_j^+Y_j^-$ occur among the
exact tests.  Consequently the equality of the pair values also holds
on every later common cone slice, by Proposition~\ref{prop:F-bridge}.

\item\label{it:G-anchor}
Either all $y_j$ vanish on $\mathcal Q$, or their vector lies on a single
rational line.  In the latter case one accumulated variable, indexed by
$0$, is an \emph{anchor}.  It was chosen at an unbounded scale $L_1$;
its pair has cost $O(L_1)$, and $L_1\le L_-$.  For each other accumulated
variable there are integers $t_j,d_j$, with $d_j>0$ bounded, such that
\begin{equation}\label{eq:G-proportions}
 y_j=(t_j/d_j)y_0,
 \qquad |t_j|=O(L_-/L_1).
\end{equation}
These equations include the point where all values are zero and are
imposed by tests of both signs.  In the all-zero case there is no anchor.

\item\label{it:G-baseline}
The optimization height $\mathfrak B$ is rational and positive on
$\mathcal Q$, and in the normalization $q(h)=1$ satisfies
\begin{equation}\label{eq:G-baseline-bounds}
                  \rho\le\mathfrak B(q)\le C
\end{equation}
for fixed positive constants.  Initially $\mathfrak B(q)=q(h)$.
It can change once, when an anchor is first chosen, to
\begin{equation}\label{eq:G-baseline}
 \mathfrak B(q)=q(h)-\mu q(Y_0^{\mathrm{sign}})/L_1,
 \qquad \mu\ge0,
\end{equation}
where either member of the anchor pair can be used and $\mu$ is bounded.
This is an allowed height in Theorem~\ref{thm:E-exact-minimizer}.  On a common
central slice the same expression has a uniform upper bound after
normalizing $q'(h)=1$, because both anchor ratios have cost $O(L_1)$.

\item\label{it:G-working}
Take the exact $\mathfrak B$-minimizer and apply the rank refinement of
Proposition~\ref{prop:F-refinement}, choosing the approximation error to
tend to zero.  This produces an equivariant centred working valuation
$w$, with finitely generated normal graded algebra, whose values on all
regular functions are multiplicatively $1+o(1)$ times those of the exact
minimizer.  Its $\mathfrak B$-normalized volume is within $1+o(1)$ of the
minimum.  Its cone satisfies the balance estimate
\eqref{eq:F-balance} and the order condition \eqref{eq:C-alpha}, with error $\eta_i\to0$, on the angular tests
$\mathcal L$ specified by the current initial tests.  The symbols of all
accumulated $Y_j^+$ are algebraically independent.  If there is an anchor,
\begin{equation}\label{eq:G-anchor-size}
                        |w(Y_0^+)|\ge cL_1
\end{equation}
for a fixed $c>0$.
\end{enumerate}

The balance hypothesis needed for the refinement in
\ref{it:G-working} follows from the preceding height bounds.
Let $v$ be the exact $\mathfrak B$-minimizer, normalized by $v(h)=1$.
For every common central $q^0$ with $q^0(h)=1$,
Corollary~\ref{cor:E-domination} gives
\[
 q^0(s)\le 2^n\frac{\mathfrak B^0(q^0)}{\mathfrak B(v)}v^0(s)
          \le\frac{2^nC}{\rho}v^0(s)
 \qquad(0\ne s\text{ regular and homogeneous}).
\]
Here only the upper bound for $\mathfrak B^0$ is used on the central
slice; the denominator is bounded below at the actual minimizer.
Proposition~\ref{prop:F-refinement} then preserves balance with a fixed
constant.

Initially there are no tests and no pairs, so these conditions follow
from the centred point $T$ and Sections~\ref{sec:E}--\ref{sec:F}.
Corollary~\ref{cor:E-height-near-minimizers} applies to
$\mathfrak B$.  Indeed, normalize two valuations by $\mathfrak B=1$.
The radially scaled version of Lemma~\ref{lem:E-convexity} fixes
the averaged value of $h=A$ and gives exact interpolation of the opposite
anchor pair.  In particular its allowed height is at most $1$.
Its volume is therefore at least the minimum in height-one normalization.
The same gap argument gives multiplicative comparison on all functions.
The bounds \eqref{eq:G-baseline-bounds} then convert that comparison back
to $h=1$.  Applied to numerators and $h$-powers, it controls ratio values
up to the corresponding cost times the relative error.

Let $M,\mathbf u,\sigma$ be the invariant character lattice, the weight of
$h$, and the cone from Proposition~\ref{prop:D-cone-dictionary}, for the
working algebra.  Equip $M_\R$ with
\begin{equation}\label{eq:G-width}
 \|k\|=\max_{\substack{y\in\sigma\\\langle y,\mathbf u\rangle=1}}
                         |\langle y,k\rangle|.
\end{equation}
Fine weights evaluate injectively under $w$.  Thus the exact equations
in \eqref{eq:G-proportions} imply the same equations for the weights of
the symbols.  In the line case write
$\kappa_0=\operatorname{wt}(\operatorname{in}_wY_0^+)$.  Opposition on pure
common tests and the paired costs give $\|\kappa_0\|=O(L_1)$, whereas
\eqref{eq:G-anchor-size} gives the reverse bound.  Hence
\begin{equation}\label{eq:G-anchor-width}
                         \|\kappa_0\|\asymp L_1.
\end{equation}

After a stage with $L_-\to\infty$, we will additionally retain:
\begin{enumerate}[label=\textup{(H\arabic*)}]
\item\label{it:G-lattice}
Every lattice vector of width $O(L_-)$ belongs to $\Q\kappa_0$ if an
anchor is present, and is zero otherwise.  In the line case the index of
$\Z\kappa_0$ in $M\cap\Q\kappa_0$ is bounded.
\item\label{it:G-field}
The angular tests in $\mathcal L$, namely $a=0$ with nonnegative orders
on all initial tests, restrict trivially to the field generated by
$S(L_-D)$.
\end{enumerate}
Here $a,D,S(\cdot)$ have the angular meanings of
Sections~\ref{sec:C}--\ref{sec:D}.  No assertion
\ref{it:G-lattice} or~\ref{it:G-field} is required while $L_-$ remains
bounded, and there is then no anchor.  These two assertions will be proved
for the next stage after the high-scale pin, without being used to
construct its rank refinement.

\subsection{The next character block and the terminal case}

Suppose first that an anchor is present and $M$ has no direction
transverse to its line.  Write $\kappa_0=\ell k_0$, where $k_0$ is
primitive in $M$; assertion~\ref{it:G-lattice} bounds $\ell$.  The
restriction of $w$ to the actual quotient field has rational rank one
and is Abhyankar.  To justify that this is the actual value group, take
invariants in every level of the finite-group-stable filtration:
exactness of invariants identifies the invariant graded algebra with the
graded algebra of the quotient filtration.  Its homogeneous fraction
weights generate $M$, and evaluation is injective.  Thus the quotient
value group is generated by
\[
                         \gamma=|w(Y_0^+)|/\ell.
\]
Abhyankar equality is preserved under the finite field extension, so this
rank-one valuation is divisorial up to scaling.  Its primitive
normalization is $w|_{\Frac(S/\Gamma)}/\gamma$.  Quasi-\'etale log pullback
and $A_S(w)=1$ give its discrepancy as $1/\gamma$.  By
\eqref{eq:G-anchor-size} and $L_1\to\infty$, this tends to zero,
contradicting $\epsilon$-lc of the actual quotient.  Rank zero of $M$ is
impossible, since a finite extension of the quotient field cannot support
a nontrivial valuation with trivial restriction.  This disposes of the
terminal possibilities.

Otherwise choose a block of lattice directions as follows.  With an
anchor use the quotient norm modulo $\R\kappa_0$ on the projected
lattice; without one use \eqref{eq:G-width}.  After a previous high stage
its shortest nonzero length divided by $L_-$ tends to infinity.  For if a
projected vector had length $O(L_-)$, lift it and subtract the nearest
multiple of the primitive line vector.  Its remaining line contribution
is $O(L_1)$ by \eqref{eq:G-anchor-width} and the bounded index, producing
a transverse lattice vector of width $O(L_-)$, contrary to
assertion~\ref{it:G-lattice}.

Before the first high stage, if some successive minima remain bounded,
choose a nonempty maximal bounded block of independent directions,
of width $O(L)$ with $1\le L=O(1)$.  This is a \emph{bounded stage}.
In every other case choose an integral scale $L\to\infty$ comparable to
the first successive minimum, and take the whole first block of
successive minima comparable to $L$.  After a subsequence, the next
successive minimum, if any, divided by $L$ tends to infinity.  These are
\emph{high stages}.  Lift the selected directions $k_j$ to $M$ by rounding
along the anchor line when necessary.  They have width $O(L)$, and
$L_-=O(L)$.

In a high stage, let $s$ denote the block size plus one when an old anchor
is present, and just the block size otherwise.  Put
\begin{equation}\label{eq:G-Rstar}
 R_* = L/L_1\quad\text{with an old anchor},
 \qquad R_*=1\quad\text{without one}.
\end{equation}
For every fixed $N\ge1$, eventually
\begin{equation}\label{eq:G-lattice-count}
 \#\{k\in M:\|k\|\le CNL\}\le C'R_*N^s,
\end{equation}
with $C'$ independent of $N$.  Indeed, the gap to the next successive
minimum confines these vectors to the first-block span modulo the line.
Distinct projected lattice points have separation at least $cL$, so a
norm-ball packing in its bounded dimension gives $O(N^{s-1})$ points
when there is an anchor, and $O(N^s)$ otherwise.  Above each projected
point, the number of lifts in the ball is $O(NL/L_1)$, since the lattice
step on the anchor line has width comparable to $L_1$.  This proves
\eqref{eq:G-lattice-count}.

\subsection{Neutral tests and simultaneous non-cancellation}

Apply Lemma~\ref{lem:C-saturation-count} at scale $L$.
Its test-cost hypothesis uses $L_-=O(L)$; its additional field hypothesis,
when $L_-\to\infty$, is precisely assertion~\ref{it:G-field}.
Write its outputs as $K_*,p,v_i,B$.  If the angular field is the constant
field, set $p=0$ and $v_i=1$.  Increase the fixed constant $B$ if needed.
Lift a basis of $S(BLD)$ to actual ratios $G_a$ of cost $O(L)$ with the
prescribed neutral initial symbols.  This basis is finite for each germ,
but its cardinality may grow along the sequence; the cost constant is
uniform.  We include $1$ in the basis.  We may
also arrange that $p$ basis members have algebraically independent
initial symbols: the cutoff plateau has an endpoint $\beta_1\ge c/L$,
and its bounded birational system in $D_W/\beta_1$ pulls back into
$S(BLD)$ and generates $K_*$ as a field.  A transcendence basis can be
chosen from a field-generating set and extended to a vector-space basis.

For each selected direction $k_j$, Corollary~\ref{cor:D-regularization}
supplies homogeneous invariant ratios of weights $\pm q_jk_j$, with
bounded positive integers $q_j$, and cost $O(L)$.  Align the two
multipliers by taking bounded powers and lift the ratios to pairs
$Y_j^\pm$ on the actual germ.  Its small-height hypothesis holds: on a
primitive angular test with $a=0$ and $LD\to0$, an old test of cost
$O(L_-)$ has negative order bounded in absolute value by $O(L_-D)=o(1)$.
Its order is integral, hence is nonnegative eventually.  Such a test
therefore belongs to $\mathcal L$, where the required
\eqref{eq:C-alpha} holds.

Add the $G_a$ and the new products $Y_j^+Y_j^-$ as tests, before imposing
any individual new pair or proportionality constraint.  Call the result
the \emph{pre-pinning slice}.  All these tests have value zero at $w$,
so this slice contains $w$.  On the working cone, the old tests together
with the $G_a$ tests force trivial restriction on $K_*$ by saturation:
nonnegativity on a vector-space basis implies nonnegativity on its whole
linear system.  The assertion also holds for pure tests, which are
trivial on the angular field.  Lemma~\ref{lem:D-lc-tests} supplies the
homogeneous Gauss description and the finite-cover lifts for all other
common lc tests.  Rational monomial tests suffice, and the equalities
extend to their real subcones.  The $G_a$ and the new pair products have
initials in $K_*$, since their neutral costs are $O(L)$.  The bridge
therefore gives value at most zero on the actual pre-pinning slice.
Their test inequalities give the reverse sign, proving exactness and
opposition there.

\begin{lemma}[Non-cancellation in every degree]\label{lem:G-noncancellation}
At the preceding stage, consider any finite linear combination of
monomials in the accumulated and newly added $Y_j^+$ and the current
$G_a$, with nonnegative exponents.  Suppose its least-$w$ group has a
nonzero sum of initial symbols.  If all its terms have the same value at
an actual point of the pre-pinning slice, their sum has that value.
This assertion holds eventually simultaneously in all polynomial degrees.
Moreover, at a later working valuation on this slice where these terms
are tied, the same non-cancellation assertion transfers to its cone at
every common test where the terms remain tied and the pair values are
opposite.
\end{lemma}

\begin{proof}
Write $A_j=\operatorname{in}_wY_j^+$ and
$B_j=\operatorname{in}_wY_j^-$.  The pair products $A_jB_j$ lie in $K_*$.
With an old anchor, the finitely many old relations
\[
 \operatorname{wt}(A_j)=(t_j/d_j)\kappa_0
\]
give neutral rational functions $A_j^{d_j}A_0^{-t_j}$.  Each is a quotient
of two neutral functions of cost $O(L)$.  More explicitly, for $t_j\ge0$
write it as
\[
 \frac{A_j^{d_j}B_0^{t_j}}{(A_0B_0)^{t_j}};
\]
for $t_j<0$ its numerator is simply $A_j^{d_j}A_0^{|t_j|}$.  The required
costs are bounded by a fixed multiple of
$d_jL_-+|t_j|L_1=O(L_-)=O(L)$.  There are at most $n$ accumulated
variables, and their denominators $d_j$ are bounded.  Thus one fixed
constant works for this entire finite collection.  The inclusion
$S(CLD)\subset K_*$ puts their numerators and denominators in $K_*$ at
one eventual index.  Without an anchor all the old $A_j$ themselves are
neutral and of bounded cost, and belong to $K_*$ instead.

The new weights are independent modulo the old line, or modulo zero.
Consequently a monomial relation of total weight zero has zero net
exponent in each new variable.  Clearing the bounded old denominators
expresses a power of every remaining weight-kernel monomial in the
finite neutral generators just described.  Every such monomial is
therefore algebraic over $K_*$.  This is a simultaneous algebraic
statement: once the finite generators have entered $K_*$, there is no
further eventual index depending on the exponents.

Choose a term in the least-$w$ group, and multiply the whole sum by the
opposite factors $Y_j^-$ corresponding to its $Y_j^+$ factors.  Its
$G_a$ factors already have value zero.  The complemented sum has a
nonzero $w$-initial of weight zero.  Each term of that initial is
algebraic over $K_*$ by the preceding kernel argument, so the initial
itself is algebraic over $K_*$.  Every common cone test is trivial on
this element.  Here we use the elementary fact that a valuation trivial
on a field is trivial on its algebraic extension inside a valued
field: a polynomial equation with nonzero constant term would otherwise
have a unique term of smallest value.

The complemented sum is an invariant ratio with a regular numerator
after multiplication by some power of $h$.  Its cost is allowed to
depend on the chosen polynomial.  Proposition~\ref{prop:F-bridge} has
no uniform cost restriction, and yields an actual upper bound zero on
the pre-pinning slice.  At a point where the original terms tie,
opposition and the zero values of $G_a$ give value zero for every
complemented term.  The valuation inequality gives the lower bound
zero.  The complemented sum thus has exact value zero, which proves
non-cancellation.

At a later working valuation $w'$ where the terms tie, this establishes
exact value zero before any new bridge is invoked.  Its initial symbol
there is nonzero.  Apply the bridge in the other direction, using only
tests exact at $w'$.  At a common test on the new cone where the terms
remain tied, the transferred upper bound and the same termwise lower
bound again give equality.  This proves the onward assertion.
\end{proof}

The enlarged collection of positive symbols is algebraically independent
at $w$: the new weights are independent modulo the old line, and the
old symbols were independent.  It remains independent at every later
working valuation on the pre-pinning slice.  Indeed, a proposed
polynomial relation among the later symbols can be grouped by their
later grading.  Each nonempty group has a nonzero least-part initial at
$w$, and Lemma~\ref{lem:G-noncancellation} forbids cancellation of that
group at the later valuation.

At a bounded stage, add both new members $Y_j^\pm$ individually as tests.
Their values must then be zero.  The original centred point $T$ remains
in the slice by \eqref{eq:G-input}, because every cost used so far is
bounded.  This invocation is uniform over the finite interval of
possible integral costs, and applies to every numerator in each of the
corresponding semi-invariant spaces.  Thus the bounded stage is complete:
replace $L_-$ by $L$ and optimize and refine again.  No assertion of
\eqref{eq:G-input} at an unbounded cost has been used.

\subsection{The count at a high scale}

The neutral tests now prevent cancellation and preserve the independent
symbols when the working valuation changes.  At an unbounded scale,
the input saturation condition does not directly control all newly
chosen costs.  We instead bound the number of symbols available at that
scale.  The high-scale pin will force enough exact relations to avoid
violating this bound.

We now work at a high stage.  For each fixed integer $N\ge1$, each integer
$0\le t\le CNL$, and the character $\chi^t$, the span of regular
homogeneous symbols of that character of grading value at most $CNL$
has dimension at most
\begin{equation}\label{eq:G-count}
                            C'R_*v_iN^{p+s}.
\end{equation}
The constant $C'$ is independent of $N$ and $t$.  To see this, divide a
symbol by $h^t$.  It becomes an invariant homogeneous fraction with
weight in $M$.  Balance \eqref{eq:F-balance}, tested on rational pure
rays, bounds its width by $O(NL)$.  Those rays are common tests because
all current test weights are zero.  For each weight,
Corollary~\ref{cor:D-character-count} bounds its dimension by
$C''v_iN^p$.  Sum over \eqref{eq:G-lattice-count}.

The same upper bound, with only a fixed change of constant, holds on the
cone of any other actual equivariant centred valuation $v$ with
\begin{equation}\label{eq:G-domination}
                           v\ge c w
\end{equation}
on all regular functions, in $h=1$ normalization, for fixed $c>0$.
Indeed the order ideals satisfy
$\mathfrak a_{a/c}(w)\subseteq\mathfrak a_a(v)$.  The resulting surjection
of finite colength quotients is equivariant.  Taking each finite-group
eigenspace is exact, so the cumulative dimensions of its associated
graded pieces have the asserted comparison.  The strict or weak endpoint
of a grading cutoff can be handled by slightly increasing its fixed
constant.  In particular we transfer \eqref{eq:G-count} from the original
working cone; we never replace it by a count with a constant depending
on a later tilt.

We record the elementary product estimate used below.
\begin{lemma}[Uniform product growth]\label{lem:G-product-growth}
Let $V$ be a finite-dimensional subspace of a function field over an
algebraically closed field, containing $1$.  Put $g=\dim V$ and
$a=\trdeg k(V)\le n$.  If $V^{(N)}$ is the span of $N$-fold products of
members of $V$, there is $c_n>0$, depending only on $n$, such that
\[
                       \dim V^{(N)}\ge c_n gN^a
                       \qquad(N\ge1).
\]
\end{lemma}
\begin{proof}
The rational map defined by $V$ has a nondegenerate integral projective
image of dimension $a$ in $\mathbf P^{g-1}$.  Its degree-$N$ homogeneous
coordinate space is $V^{(N)}$.  Successively cut by general hyperplanes.
In positive dimension the reduced hyperplane section is nondegenerate
in its hyperplane.  One way to see the required nondegeneracy is to pass
to the normalization: a general hyperplane divisor is reduced in
codimension one, so another linear form vanishing on its reduced support
has a regular ratio with the cutting form, and the ratio is constant.
The reduced section is integral until the last curve section by Bertini;
the final reduced set of points spans its hyperplane and has Hilbert
function at least its dimension in degree one in every higher degree.
For the homogeneous coordinate ring, multiplication by a nonzero
cutting linear form is injective, and the quotient surjects onto the
coordinate ring of the reduced cut.  Thus successive differences of
Hilbert functions dominate the Hilbert functions of these cuts.
Summing these inequalities gives, for $a>0$, the sufficient bound
\[
        \dim V^{(N)}\ge(g-a)\binom{N+a-1}{a}.
\]
Since $g\ge a+1$, one has $g-a\ge g/(a+1)$, giving the asserted constant.
For $a=0$ the image is a point, $g=1$, and the assertion is immediate.
\end{proof}

\subsection{The high-scale pin}

At a high stage the new costs grow with $L$, so the original input
\eqref{eq:G-input} no longer supplies a centred point after individual
values are forced to zero.  We instead constrain their proportions while
retaining a centred point close enough to the current minimum.

\begin{proposition}[High-scale pin]\label{prop:G-high-pin}
Consider a high stage with pre-pinning slice, working valuation $w$,
baseline height $\mathfrak B$, scale $L$, and $s$ selected values as above.
After passage to a subsequence, finitely many invariant tests of cost
$O(L)$ can be added with the following properties.
\begin{enumerate}[label=\textup{(\roman*)}]
\item The resulting slice contains a centred valuation.  The added tests
are exact on that slice, all old tests and proportionalities remain in
force, and the accumulated positive symbols remain algebraically
independent at the exact minimizer and working valuation in
\textup{(iii)} below.
\item The selected values either vanish or lie on a single rational
line.  In the line case their proportions have the bounds
\eqref{eq:G-proportions}, with $L_-=L$.  An old anchor remains the anchor;
if a new anchor is chosen, its scale is $L_1=L$.
\item The baseline height is retained unless an anchor is first chosen,
in which case it may change to the rational height
\eqref{eq:G-baseline}.  It satisfies \eqref{eq:G-baseline-bounds} on the
new actual slice.  At its exact minimizer $v'$, and at an arbitrarily close
rank refinement $w'$, the anchor, when present, satisfies
$|v'(Y_0^+)|,|w'(Y_0^+)|\ge cL_1$.  Both valuations dominate a fixed positive
multiple of $w$ on every regular function.  In particular,
\[
                         w'\ge c w.
\]
\end{enumerate}
Consequently \ref{it:G-tests}--\ref{it:G-working} hold with $L_-=L$.
The next-scale assertions \ref{it:G-lattice}--\ref{it:G-field} are not
assumed in constructing the pin or its rank refinement.
\end{proposition}

Domination preserves the counting bound on the original working cone.
The rational proportions permit equations of cost $O(L)$, and the
nonzero anchor size will detect a small primitive quotient discrepancy
if only its line remains.  We first construct this pin, then prove that
it also gives the next-scale field and lattice assertions.

\subsection{Finding all necessary relations}\label{subsec:G-pin-rows}

We begin the proof of Proposition~\ref{prop:G-high-pin}.  Within this
stage, progress is measured by the rank of the integral relations imposed
on the $s$ selected values.  This row rank is separate from the
transcendence degree used to terminate the outer induction.

For the $s$ selected variables put $L_j=L_1$ for an old anchor, and
$L_j=L$ for every new variable, and set
\[
                     x(q)=(y_j(q)/L_j)_{j=1}^s.
\]
When there is an anchor we include its index $0$ in this notation.
Opposition and the paired costs bound all coordinates uniformly on the
actual pre-pinning slice.  Lemma~\ref{lem:E-convexity} and Corollary~\ref{cor:E-opposite-ratios}
give exact convex interpolation of the coordinates, so their image is compact and
convex.  Its points can be realized by centred valuations whenever the
interpolation assigns positive weight to a centred input.

We successively impose integral rows
\begin{equation}\label{eq:G-row}
 \sum_j k_jy_j=0,
 \qquad \widehat k=(k_jL_j/L)_j,
 \qquad \|\widehat k\|_{\mathrm{Eucl}}=O(1).
\end{equation}
A row is imposed by testing two opposite monomials.  For its positive
sign use
\[
 M_k=\prod_{k_j\ge0}(Y_j^+)^{k_j}
       \prod_{k_j<0}(Y_j^-)^{-k_j};
\]
exchange the signs to obtain the other monomial.
Each has cost $O(\sum_j|k_j|L_j)=O(L)$, and opposition makes the two
inequalities equivalent to the row equality.

Let $\mathcal R_i\subset\R^s$ denote the span of the normalized rows
already imposed.  We retain a comparator $v_0$ on the current slice
whose $\mathfrak B$-normalized volume is within $1+o(1)$ of the minimum
on the actual slice at the start of this high stage.  Initially use $w$.
The gap estimate then gives $v_0=(1+o(1))w$ on all regular functions.
With an old anchor, $|y_0(v_0)|/L_1$ stays bounded away from zero.
Pass to subsequences so the row spaces converge and the limits of the
chosen rows are independent.

For clarity, independence of their transverse integer rows is equivalent
to independence of these normalized row limits in the present situation.
The non-anchor coefficients $k_j$ are bounded integers, hence become
fixed on a subsequence.  If a rational combination of their transverse
rows vanished, the same combination of full rows would be supported
only on the anchor.  Its exact value at $v_0$ must vanish.  Since the
anchor coordinate is nonzero, its anchor coefficient vanishes as well.
Thus any dependence in the transverse rows is a dependence of the full
rows.  Conversely independent transverse rows keep the full normalized
rows independent in the limit.  With no anchor, all coefficients are
bounded integers and the statement is immediate.

Let $K_i$ be the projection image $\{x(q)\}$ of the current constrained
slice, and let
\[
                   C_i=\operatorname{conv}(\{0\}\cup K_i)
                       =[0,1]K_i.
\]
After a subsequence these bounded compact convex sets converge in
Hausdorff distance.  Their limit lies in
$\mathcal R_\infty^\perp=\lim\mathcal R_i^\perp$.
We first handle the case
\begin{equation}\label{eq:G-thin}
                      \operatorname{span}(\lim C_i)
                      \ne\mathcal R_\infty^\perp.
\end{equation}
Choose a unit vector $u\in\mathcal R_\infty^\perp$ perpendicular to the
left-hand span, and rational vectors $u_i\to u$.  Then
\begin{equation}\label{eq:G-projection-small}
                    u_i\cdot x(q)=o(1)
\end{equation}
uniformly on the current slice.  Fix a sufficiently small
$\delta>0$, with its size chosen below, and consider the rational height
\begin{equation}\label{eq:G-thin-height}
 b(q)=\delta\mathfrak B(q)-\sum_j u_{ij}y_j(q)/L_j.
\end{equation}
We take the scales integral, so all coefficients can be rational.
According to the sign of $u_{ij}$, use $Y_j^+$ or $Y_j^-$, and absorb its
$h$-denominator into the coefficient of $q(h)$.  This puts $b$ in the
allowed form with nonnegative coefficients on the subtracted regular
numerator values.  Equations \eqref{eq:G-baseline-bounds} and
\eqref{eq:G-projection-small} show that it is positive and equals
$(1+o(1))\delta\mathfrak B$ uniformly on this slice.

Let $v$ be its exact minimizer.  Comparison with $v_0$ shows that the
$\mathfrak B$-normalized volume of $v$ is within $1+o(1)$ of the original
stage minimum.  Consequently
\begin{equation}\label{eq:G-thin-close}
                           v=(1+o(1))w
\end{equation}
on all regular functions, and $v$ is exact on all present tests.
On its cone every nonzero pure test, and every Gauss test over a primitive
angular divisor with $a=0$ and $LD\to0$, is common.  Indeed all test
weights are zero, their costs are $O(L)$, and the integral angular orders
cannot be negative of size $o(1)$.  All pair equalities transfer to this
cone.  Positivity of the induced height, supplied by
Theorem~\ref{thm:E-exact-minimizer}, and the uniform central upper bound for $\mathfrak B$ imply
\begin{equation}\label{eq:G-central-tilt}
                \sum_j u_{ij}y_j'/L_j\le C\delta q'(h)
\end{equation}
on these tests.

Apply Corollary~\ref{cor:D-tilted-characters}.  For a fixed
$M_0\gg1/\delta$, choose integers
\begin{equation}\label{eq:G-rounding}
                r_j=M_0u_{ij}L/L_j+O(1).
\end{equation}
It gives a homogeneous invariant ratio $Z$ on the cone of $v$, of weight
\begin{equation}\label{eq:G-cheap-weight}
 -q\sum_jr_j\operatorname{wt}(\operatorname{in}_vY_j^+),
\end{equation}
where $q$ is a bounded positive integer, and of cost at most
\begin{equation}\label{eq:G-cheap-cost}
                              CqM_0\delta L.
\end{equation}
The constant $C$ is fixed before $\delta$ and $M_0$ are chosen.  The bound
on $q$ may depend on these fixed choices.  This distinction follows from
the proof of the tilted-character result: its height correction has
size $O(M_0\delta L)$; the $O(L)$ rounding error and the fixed constant in
\eqref{eq:G-central-tilt} are absorbed by first taking
$M_0\gg1/\delta$.  Only the later regularizing multiplier need depend
on these parameters.  The absolute grading value of $Z$ has the same
bound.  In fact \eqref{eq:G-cheap-weight} evaluates to
\[
 -qM_0L\,u_i\cdot x(v)+O(qL),
\]
which is bounded by \eqref{eq:G-cheap-cost} eventually, by
\eqref{eq:G-projection-small} and the chosen size of $M_0$.
No assumption \eqref{eq:C-alpha} on this newly tilted cone is used.

\begin{claim}\label{clm:G-new-row}
After a subsequence there is an additional row of the form
\eqref{eq:G-row}, exact at $v$, whose transverse part is independent of
the previously imposed transverse rows.
\end{claim}
\begin{proof}
Suppose otherwise.  Then, for each fixed bound on normalized coefficient
size, every exact integral tie at $v$ eventually belongs to the old row
span.  To verify this assertion, a tie with transverse part outside the
old span would give the desired row after fixing its bounded integer
transverse coefficients.  If its transverse part lies in the old span,
subtract the corresponding rational combination of the fixed old
transverse rows.  The remainder is a pure-anchor tie.  It is zero
because \eqref{eq:G-thin-close} preserves the nonzero anchor value at
$v$.  Without an anchor there is no remainder.  The claimed eventual
assertion for all ties with a given bound follows by choosing a failing
tie along a subsequence if necessary.

Let $\widetilde u_i$ be a unit vector in $\mathcal R_i^\perp$ tending to
$u$.  Choose a large fixed integer $N\gg qM_0$.  In the exponent box
\[
           0\le e_j\le\lfloor NL/L_j\rfloor,
           \qquad \widehat e=(e_jL_j/L)_j,
\]
there are at least $cR_*N^s$ integer points.  The range of
$\widetilde u_i\cdot\widehat e$ has length $O(N)$.  Partition this range
into intervals of width $qM_0/4$.  One slab contains at least
\begin{equation}\label{eq:G-slab-size}
                           cR_*N^{s-1}qM_0
\end{equation}
points.  For each exponent in this slab, consider on
$\operatorname{gr}_v$ its positive-variable monomial multiplied by $Z^l$,
where
\begin{equation}\label{eq:G-powers}
                     0\le l\le\left\lfloor
                                     N/(qM_0\delta)\right\rfloor,
\end{equation}
and by a basis of the products of $N$ of the $G_a$ whose $w$-initials
are independent.  Saturation and counting supply at least $v_iN^p$
such products.

Different powers $l$ have disjoint grading values in this construction.
A collision would give the exact integral tie
\[
                              e-e'-q(l-l')r.
\]
Its normalized size is bounded with $N,M_0,\delta,q$ fixed.  Under our
supposition it belongs to the old row span, so its normalized dot
product with $\widetilde u_i$ is zero.  The slab bounds the contribution
from $e-e'$ by $qM_0/4$.  But \eqref{eq:G-rounding} gives
\[
       \widehat r=M_0u_i+O(1),
       \qquad |\widetilde u_i\cdot\widehat r|\ge M_0/2
\]
eventually, after increasing the fixed $M_0$.  For $l\ne l'$ the other
contribution has magnitude at least $qM_0/2$, a contradiction.  It is
essential here that $\widetilde u_i$ is orthogonal to the entire old
normalized row span.

At any one power $l$, the slab size times $v_iN^p$ symbols are linearly
independent.  Before changing to $v$, the selected variable weights
separate the exponent vectors, and the chosen $G_a$ products have
independent initials at $w$.  Any possible relation after the change
can be grouped by its $v$-grading, and
Lemma~\ref{lem:G-noncancellation} rules out cancellation in each group.
Multiplication by the nonzero $Z^l$ preserves this independence.
Combining with disjoint grading values for different $l$ gives at least
\begin{equation}\label{eq:G-excess-count}
                     c\delta^{-1}R_*v_iN^{p+s}
\end{equation}
independent ratios.

Every ratio just counted has cost and absolute grading value at most
$C_0NL$, where $C_0$ is independent of $\delta,M_0,q$.  The slab
monomials cost $O(NL)$; the $G_a$ products do too; and
\eqref{eq:G-cheap-cost} multiplied by the exponent bound
\eqref{eq:G-powers} costs $O(NL)$.  A common $h$-denominator gives
regular numerators of one finite character and degree $O(NL)$.
By \eqref{eq:G-thin-close}, eventually $v\ge w/2$, so the transferred
upper bound \eqref{eq:G-count} has a constant independent of the tilt.
It contradicts \eqref{eq:G-excess-count} if $\delta$ is sufficiently
small.

The order of choices is as follows.  First fix the current stage
constants, the old rows, the saturation constant and the original
counting constant.  Then fix $\delta>0$ small enough for this last
contradiction.  Next fix $M_0$ sufficiently large compared with
$1/\delta$ and the rounding constants.  The character construction then
gives a bound $Q$ for $q$; pass to a subsequence fixing $q$, or choose
using $Q$.  Only now choose a sufficiently large fixed $N$, and finally
let $i$ be large enough for all of these fixed-parameter assertions.
Thus $\delta$ never tends to zero with $i$, and no eventual count is
used at a variable product degree.  This proves the claim.
\end{proof}

Impose the new row by both signs and retain $v$ as the new comparator.
It is centred and still within $1+o(1)$ of the original stage minimum.
The limits of the enlarged rows stay independent because their
transverse integer projections are fixed and independent.  Each time
\eqref{eq:G-thin} holds, Claim~\ref{clm:G-new-row} increases this rank.
There are at most $s$ such increases, and at most $s-1$ with a retained
nonzero anchor.  Therefore this procedure reaches the alternative
\begin{equation}\label{eq:G-full-span}
                    \operatorname{span}(\lim C_i)
                    =\mathcal R_\infty^\perp.
\end{equation}
This is strict finite progress within a stage.

\subsection{Rational directions and anchor size}\label{subsec:G-pin-directions}

Fix a small positive number $\lambda$.  Mix the comparator projection
with the full current projection set:
\[
                  K_i^\lambda=(1-\lambda)x(v_0)+\lambda K_i.
\]
These points have actual centred realizations $q$ with
$q\ge(1-\lambda)v_0$ on all regular functions, by Lemma~\ref{lem:E-convexity}.  Their baseline heights interpolate exactly because the
height depends only on $h$ and the anchor coordinate.  The bounds on
$\mathfrak B$ show, for sufficiently small fixed $\lambda$, that their
$\mathfrak B$-normalized volumes are within any prescribed small fixed
relative error of the original stage minimum.  For example domination
gives the volume factor $(1-\lambda)^{-n}$, and the height changes by
$O(\lambda)$ relative to its positive lower bound.

The radial hull $[0,1]K_i^\lambda$ still has a full-dimensional limit in
$\mathcal R_\infty^\perp$.  Indeed $K_i^\lambda$ contains $x(v_0)$,
and differences from that point span
$\lambda(K_i-x(v_0))$.  Its linear span is therefore the same as that of
$K_i$.  The same statement holds in the limit.  A compact convex set of
full linear span containing zero has nonempty relative interior, so by
Hausdorff convergence these radial hulls contain relative balls of
radius bounded below in $\mathcal R_i^\perp$, unless that space has
dimension zero.  If there is an old anchor, choose such a ball away from
anchor-coordinate zero: the comparator has nonzero anchor coordinate,
and a convex combination with an interior point gives an interior
point retaining that property.  Radial truncation is used only to choose
a direction; every nonzero direction chosen in it has an actual mixed
representative, with the centred domination just established.

With an old anchor, write an integral row matrix as $(b_i,A)$, where
$A$ is its transverse part.  After the subsequences already taken,
$A$ is a fixed integral matrix of full row rank, and the entries of
$b_i$ are $O(R_*)$.  The feasible ratios $z_j=y_j/y_0$ satisfy
\begin{equation}\label{eq:G-rational-affine}
                                Az=-b_i.
\end{equation}
Directions in the ball away from $x_0=0$ give a relative ball of feasible
ratios, in coordinates $z/R_*$.  We can choose a point there with bounded
common denominator and $z_j=O(R_*)$.  To make the denominator assertion
explicit, a fixed invertible minor of $A$ gives a particular solution of
\eqref{eq:G-rational-affine} with denominator dividing its determinant
and size $O(R_*)$.  The rational kernel of $A$ has a fixed lattice basis.
Round in a fixed sufficiently fine denominator refinement of that
lattice.  Since $R_*\ge1$, its mesh after division by $R_*$ is uniformly
small enough to enter the specified ball.  If the affine solution
space has dimension zero, the particular solution is already the
required exact ratio.  Thus the denominator is bounded and the numerator
bounds are $O(L/L_1)$.

Without an old anchor, the full row matrix is fixed integral.  If its
orthogonal space is nonzero, the ball contains a nonzero rational
direction of bounded complexity.  Choose a nonzero coordinate as the
new anchor and put $L_1=L$.  Choose the ball and direction so that this
coordinate has absolute value at least a fixed positive constant; its
actual mixed representative has that same lower bound or a larger one,
since the radial factor lies in $(0,1]$.  If the orthogonal space has
dimension zero, all selected values are already zero.

Impose the chosen proportions, or all zero values, by both-sign
monomial tests.  Their costs are $O(L)$: with an old anchor a row
$d_jy_j-t_jy_0=0$ has cost $O(d_jL+|t_j|L_1)$ with $d_j$ bounded and
$|t_j|=O(L/L_1)$; with a new anchor both coefficients are bounded.
The selected actual mixed representative survives all these equations,
so the new slice contains a centred point.  Old proportionalities are
retained by the old tests.  In the cases with no new anchor, the exact
baseline minimum on the pinned slice is as close to the preceding
working valuation $w$ on all functions as desired, by choosing
$\lambda$ sufficiently small and applying the gap estimate on the
original stage slice.  In particular an old anchor still has
$|y_0|\ge cL_1$, and domination by a fixed multiple of $w$ holds.

When a new anchor is chosen, a further step is needed: equations on
ratios alone do not prevent the minimizing vector from moving to zero.
The prior baseline is then $q(h)$.  Let $z$ denote the centred mixed
comparator on the pinned slice.  Choose the sign of the anchor so that
\[
 p_0(q)=q(Y_0^{\mathrm{sign}})/L_1,
 \qquad p_0(z)\ge c_1>0.
\]
Its ordinary volume is within a prescribed small fixed error of the
prior ordinary minimum $V_*$.  Let $M=\max p_0$ on this compact,
$h=1$-normalized slice.  Thus $M\ge c_1$ and $M=O(1)$.  For
$0\le\mu\le(1-\rho)/M$, consider
\begin{equation}\label{eq:G-anchor-height}
                              b_\mu(q)=q(h)-\mu p_0(q),
\end{equation}
where $\rho>0$ is a small fixed number.  These heights are bounded below
by $\rho$ and uniformly above.  The exact minimizing ray varies
continuously with $\mu$, by Theorem~\ref{thm:E-exact-minimizer}.

At the endpoint $\mu_*=(1-\rho)/M$ the minimizer $v_*$ has
$p_0(v_*)\ge M/2$ for small enough $\rho$.  Here are the estimates.
Interpolate a maximizing witness for $p_0$ with $z$, assigning weight
$\rho$ to $z$.  The interpolant $u$ is centred, has
$b_{\mu_*}(u)\le2\rho$, and \eqref{eq:E-volume-interpolation}, applied
with $f=h$ and both inputs normalized by $h=1$, gives
\[
                         \vol(u)\le\rho^{1-n}\vol(z).
\]
Minimality and $\vol(v_*)\ge V_*$ yield
\[
 b_{\mu_*}(v_*)
 \le2\rho^{1/n}\bigl(\vol(z)/V_*\bigr)^{1/n}<\tfrac12
\]
for sufficiently small fixed $\rho$.  Therefore
$p_0(v_*)>M/[2(1-\rho)]\ge M/2$.

If the minimizer at $\mu=0$ already has $p_0\ge c_1/4$, use $\mu=0$.
Otherwise continuity gives a nonempty open interval of parameters where
$c_1/4<p_0(v_\mu)<c_1/3$, and we choose a rational $\mu$ in that
interval.  Since $p_0(z)\ge c_1$, one has
$b_\mu(z)<b_\mu(v_\mu)$.  Tilted minimality consequently gives
\begin{equation}\label{eq:G-anchor-volume}
 \vol(v_\mu)
 \le\left(\frac{b_\mu(z)}{b_\mu(v_\mu)}\right)^n\vol(z)
 \le\vol(z).
\end{equation}
The same final bound is immediate at $\mu=0$.  Thus the new minimizer
still lies as near the original ordinary minimum as required, while
its anchor value is bounded below by $c_1L_1/4$.  The gap estimate gives
the required all-functions domination.  Adopt this rational height as
$\mathfrak B$ for all subsequent stages.  It has the form \eqref{eq:G-baseline} and bounds \eqref{eq:G-baseline-bounds}.

Finally take the rank refinement with error tending to zero, sufficiently
small to preserve the anchor lower bound.  This preserves the pinning
equalities, the exact pair and test identities, and the independent
symbols by Lemma~\ref{lem:G-noncancellation}.  It supplies balance and
\eqref{eq:C-alpha} and does not use the next-scale field and lattice
assertions.  Thus \ref{it:G-tests}--\ref{it:G-working} hold with $L_-=L$, proving
Proposition~\ref{prop:G-high-pin}.  It remains to verify the two
next-scale assertions.

\subsection{Field and lattice consequences of the pin}\label{subsec:G-pin-consequences}

\begin{lemma}\label{lem:G-pin-consequences}
The working valuation $w'$ furnished by Proposition~\ref{prop:G-high-pin}
satisfies \ref{it:G-lattice}--\ref{it:G-field} at the next scale $L_-=L$.
\end{lemma}
\begin{proof}
The proposition supplies the new exact tests, rational proportions and,
when present, the anchor lower bound.  It also gives $w'\ge c w$ on all
regular functions, so the transferred count \eqref{eq:G-count} remains
available with a fixed constant.

Use double bars for the initial symbols at $w'$.  Let $E$ be the field
generated by the $p$ chosen algebraically independent
$\overline{\overline{G_a}}$ and the $s$ selected variable symbols.
These have transcendence degree $p+s$, by
Lemma~\ref{lem:G-noncancellation}: they were independent together at $w$,
since the selected variable weights were independent and the chosen
$G_a$ were neutral.  Write $E^0$ for its subfield of weight zero.
Every common test on the new cone is trivial on $E^0$.  Indeed, the
both-sign pinning equations also hold on that cone.  If there is an
anchor, its nonzero value at $w'$ ensures that two selected monomials
with equal weight at $w'$ have equal value at every common test; in the
all-zero case this is immediate.  Lemma~\ref{lem:G-noncancellation} then
gives the predicted exact value for each nonzero homogeneous polynomial.
Ratios of homogeneous polynomials of equal weight generate $E^0$,
proving the assertion.

There is also a space $V_0$ of invariant ratios on this new cone, of
cost and absolute grading value $O(L)$, with
\begin{equation}\label{eq:G-large-space}
                       \dim V_0\ge cR_*v_i.
\end{equation}
Take the double-bar symbols of the whole $G_a$ basis, and, if an old
anchor was present, multiply by its powers with exponent from $0$ to
$\lfloor R_*\rfloor$.  Their costs remain $O(L)$.  They are linearly
independent at $w$ by their distinct anchor weights and the basis
property, and remain so by Lemma~\ref{lem:G-noncancellation}.

Choose an enlargement $V$ of $V_0$, with the same type of fixed cost and
grading bounds, containing $1$, the specified generators of $E$, and every
neutral function in $S(LD_{\mathrm{next}})$.  We may also include in $V$ a
nonzero homogeneous ratio of weight a bounded positive multiple of any
chosen new lattice vector $k$ of width $O(L)$.  For the latter use
Corollary~\ref{cor:D-regularization} at $w'$: the refined cone satisfies
\eqref{eq:C-alpha}, and tiny $LD_{\mathrm{next}}$ forces all new tests
by the same integrality and cost argument as before.  The space $V$
is spanned by homogeneous ratios.  Give all products of $N$ members of
$V$ one common denominator $h^t$, where $t=O(NL)$; their regular
numerators then have a single finite character and degree $O(NL)$.
The transferred upper bound \eqref{eq:G-count} applies.

If the field generated by $V$ had transcendence degree $a>p+s$,
Lemma~\ref{lem:G-product-growth} and \eqref{eq:G-large-space} would give
at least $cR_*v_iN^a$ independent products, against
$C'R_*v_iN^{p+s}$.  Choose one sufficiently large fixed $N$ and then let
$i$ be sufficiently large.  This contradiction shows that the field
generated by $V$ is algebraic over $E$.  Uniformity for the chosen $k$ follows by
applying this argument to any alleged sequence of failing choices.

A neutral element algebraic over $E$ is algebraic over $E^0$.  In fact,
clear denominators in an algebraic equation so that its coefficients
are polynomials in the homogeneous generators of $E$, take a nonzero
homogeneous part of the equation, and divide by a nonzero coefficient of
that weight.  All resulting coefficients have weight zero.  Common
tests are trivial on $E^0$, and hence on these algebraic elements.  The
Gauss description and finite-cover lifting translate this back to every
angular test in $\mathcal L$, proving assertion~\ref{it:G-field} for the
next scale.

The same algebraic equation argument shows that the weight of an
algebraic homogeneous element lies in the rational span of the weights
of $E$.  The latter is the anchor line or zero.  This proves the first
part of assertion~\ref{it:G-lattice}.  To bound its index in the line
case write $\kappa'_0=\ell k'_0$ with $k'_0$ primitive.  Its width is
$O(L_1)$ by the paired anchor costs.  If $\ell\to\infty$, apply the
bounded-multiple construction at the current scale $L$ to all $jk'_0$,
$1\le j\le\ell$.  Their widths are $O(L_1)$ and $L_1\le L$.
There are symbols of weights $q_jjk'_0$, with
$1\le q_j\le Q$ for fixed $Q$, and cost and absolute grading value
$O(L)$.  These bounds are uniform over all these choices: otherwise a
sequence of failing $j_i$ would contradict the arbitrary-sequence
version of Corollary~\ref{cor:D-regularization}.

Let $D_0$ be a fixed common multiple of the bounded denominators of the
finitely many pinned variable weights.  The subgroup generated by those
weights has primitive step $h_i\ge\ell/D_0$, and $h_i$ divides $\ell$.
Among the integers $q_jj$, an integer has at most $Q$ preimages, and each
residue modulo $h_i$ has at most $QD_0+1$ representatives in
$[1,Q\ell]$.  Hence at least
\[
                         \frac{\ell}{Q(QD_0+1)}
\]
distinct residues occur.  All weights of the original space $V_0$ in
\eqref{eq:G-large-space} lie in this subgroup.  With an old anchor its
basis consists of neutral symbols multiplied by powers of that anchor;
when the pin has just chosen a new anchor, $V_0$ consists only of neutral
symbols.
Multiplying this original $V_0$ by symbols representing distinct residues yields a
direct sum, since the weight supports are disjoint.  Its dimension is
an unbounded multiple of $R_*v_i$, while its costs and degrees are
$O(L)$.  The fixed-$N$ version of \eqref{eq:G-count} forbids this.
Thus $\ell$ is bounded, completing assertion~\ref{it:G-lattice}.

\end{proof}

\subsection{Termination}

\begin{proof}[Proof of Theorem~\ref{thm:G-pinning}]
Assume the excluded sequence exists.  The initial data satisfy
\ref{it:G-tests}--\ref{it:G-working}.  At every stage, either the sole
remaining anchor line gives the primitive-discrepancy contradiction
above, or a nonempty independent block of new directions can be chosen.
A bounded block is pinned directly by \eqref{eq:G-input}; a high block
is pinned by the finite row procedure, the rational-direction choice,
and, if required, the anchor-height adjustment.  The next-stage
invariants then hold by the field and lattice argument.  Each
nonterminal stage adds at least one new positive variable, and their
initial symbols remain algebraically independent at the next working
valuation.  A function field of dimension $n$ cannot contain more than
$n$ algebraically independent elements.  Thus the process cannot have
more than $n$ nonterminal stages, whereas all its terminal alternatives
have been excluded.  This contradiction proves the theorem.
\end{proof}
